\section{模型解释的数学形式化}
\label{sec:14-Deep-Learning/10-Interpretability-Mathematics/04}

一个信用评分模型，输入 30 个特征——收入、负债比、近期查询次数、账户年龄等等——输出 0 到 1 之间的分数。某位客户被拒，合规部门要求书面回答：这次决策的主要原因是哪个特征。

你有三种本能。对输入求导，梯度绝对值最大的是``近期查询次数''。把``近期查询次数''替换成人群均值再跑一遍，分数从 0.32 跳到 0.61，变化最大的也是它——不过第二名从梯度法的``账户年龄''换成了``收入''。同事用 Shapley 值跑出来，第一名是``负债比''。

三份排名摆在桌上，很自然会问哪一份是对的。这个问题问不出答案，因为``哪个特征重要''还不是一个数学命题：相对哪个输出、相对哪个参照、用哪类扰动测量、在哪个分布上平均，四件事一件都没定。定下来之后，三种方法各自回答的是三个不同的问题，答案不一致是它们该有的样子。

\subsection{敏感度、必要性与充分性是三个不同的量}

把三种本能写成公式，差别立刻显形。记 \(x\) 为待解释的输入，\(b\) 为参照输入（基线），\(x_{-i}\) 表示除第 \(i\) 维之外的部分。

敏感度问的是局部斜率：
\[
s_i(x)=\frac{\partial f}{\partial x_i}(x).
\]
必要性问的是拿掉之后掉多少：
\[
n_i(x)=f(x)-f\bigl(b_i,\;x_{-i}\bigr).
\]
充分性问的是只给它能撑起多少：
\[
u_i(x)=f\bigl(x_i,\;b_{-i}\bigr)-f(b).
\]

三个量在逻辑上互不蕴含，用两个两输入的布尔函数就能看清。取 \(f=x_1\wedge x_2\)、\(x=(1,1)\)、基线 \(b=(0,0)\)：拿掉任一维输出从 1 掉到 0，两维都必要；只给任一维输出仍是 0，两维都不充分。换成 \(f=x_1\vee x_2\)，同样的 \(x\) 与 \(b\)：拿掉任一维输出还是 1，两维都不必要；只给任一维输出就是 1，两维都充分。两种情形下的 Shapley 值都是 \((0.5,0.5)\)——同一组归因分数下面，藏着两种完全相反的依赖结构。要区分它们，只能另外去问必要性和充分性。

饱和给出另一种分歧。

\begin{center}
\begin{tikzpicture}[x=1.45cm,y=3.4cm,>=stealth]
\draw[->] (-0.15,0) -- (6.4,0) node[right,font=\small] {\(x_i\)};
\draw[->] (-0.15,0) -- (-0.15,1.15) node[above,font=\small] {\(f\)};
\draw[very thick] plot[smooth] coordinates
  {(0,0.03) (0.75,0.06) (1.5,0.13) (2.25,0.28) (3,0.50)
   (3.75,0.72) (4.5,0.87) (5.25,0.945) (6,0.975)};
\fill (0.4,0.042) circle (1.8pt);
\node[below,font=\scriptsize] at (0.4,0.02) {\(b_i\)};
\fill (5.25,0.945) circle (1.8pt);
\node[below right,font=\scriptsize] at (5.3,0.93) {\(x_i\)};
\draw[very thick,dashed] (4.35,0.905) -- (6.15,0.995);
\node[above,font=\scriptsize] at (4.9,0.99) {切线几乎水平};
\draw[dotted] (0.4,0.042) -- (6.25,0.042);
\draw[dotted] (5.25,0.945) -- (6.25,0.945);
\draw[<->] (6.25,0.042) -- (6.25,0.945);
\node[right,font=\scriptsize,align=left] at (6.3,0.5) {路径上累积的\\ 贡献 \(\approx0.90\)};
\end{tikzpicture}
\end{center}

\emph{同一个特征，在当前点的梯度接近零，从基线走到当前点的累积贡献却接近满分。}

图里那条曲线在 \(x_i\) 处已经进入平台，切线几乎水平，梯度归因据此报告这个特征不重要。而从基线 \(b_i\) 走到 \(x_i\) 的过程中，输出上升了 0.90，路径类归因据此报告它贡献极大。两个数字都算对了，只是一个测的是此刻再动一点会怎样，另一个测的是一路走来发生了什么。合规场景要的多半是后者，而模型对输入噪声的鲁棒性分析要的是前者。

\subsection{Shapley 值把分配公理化，代价是必须发明``缺失''}

Shapley 值处理的是另一个层次的问题：怎样把 \(f(x)-f(b)\) 这个总差额公平地分给各个特征。给定集合函数 \(v(S)\)——只使用特征子集 \(S\) 时的输出——特征 \(i\) 的分配额是它在所有加入顺序下边际贡献的加权平均：
\[
\phi_i=\sum_{S\subseteq N\setminus\{i\}}
\frac{\abs{S}!\,(n-\abs{S}-1)!}{n!}
\bigl[v(S\cup\{i\})-v(S)\bigr].
\]
那个组合权重是在全部 \(n!\) 种引入顺序中，恰好以 \(S\) 为 \(i\) 的前驱集合的顺序占比。效率、对称性、虚拟参与者、可加性四条公理唯一确定了这个分配方案，这是它被广泛采用的理由。

难点不在公式，在 \(v(S)\)。模型要求全部 30 个输入才能前向传播，只给 \(S\) 时缺的那些得填。用边缘分布采样，问的是``在全体人群中，已知 \(S\) 上的取值、其余按总体分布看''；用条件分布采样，问的是``在 \(S\) 取这些值的那部分人群中，其余通常取什么''；用一个具体的背景样本填充，问的是``以这位参照人为基准''。三种填法对应三个不同的 \(v\)，因而对应三组不同的 \(\phi\)。

特征独立时前两种等价，线性模型下还化简为 \(w_i(x_i-\E[X_i])\)，看起来很像人们想要的东西。特征相关时——收入与负债比通常负相关——共享信息该记在谁头上，完全由填充策略决定。归因结果因此不是模型的属性，而是模型、特征分布与缺失处理三者合成的属性。报告 Shapley 值却不报告 \(v\) 的定义，等于报告了一个无法复现的数。另外，精确计算需要遍历 \(2^n\) 个子集，实践中一律用采样或结构假设近似，近似误差同样需要报告。

\subsection{LIME 解释的是替身，不是模型}

LIME 换了一条路：不追求公理，只在目标样本邻域里用一个人类读得懂的简单模型去逼近原模型。做法是在 \(x\) 附近按某种扰动机制生成一批变体——随机遮蔽若干词、用背景值替换若干像素——用原模型给它们打分，再做距离加权的稀疏线性回归，把系数报告为重要性。

它的陈述范围应当写全：在选定的邻域形状、扰动分布、核宽和替身模型类下，这个简单函数最接近原模型在该区域的行为。三个自由度都会改变结论。邻域太小，线性替身退化成局部梯度；太大，替身根本拟合不住，系数不再有意义。扰动机制若生成的是数据分布之外的样本——对图像做大块矩形遮挡产生的东西不像任何真实图像——那么替身解释的是模型在人造输入上的行为。替身模型类选线性只能捕一阶依赖，选树能捕交互却失去系数排序的直观。

模型足够光滑、邻域足够小、采样覆盖充分时，LIME 的系数趋向局部梯度，与梯度归因收敛到同一个答案。一般情形下两者不必一致，所以报告 LIME 时应当附上邻域定义、替身的拟合误差，以及多次重跑的稳定性。

\subsection{积分梯度用路径买到完备性，用基线付账}

积分梯度从一条要求出发：输入与基线之间的输出差额要被完全分配，不留残差。沿直线路径累积梯度，
\[
\mathrm{IG}_i(x)=(x_i-b_i)\int_0^1\frac{\partial f}{\partial x_i}\bigl(b+\alpha(x-b)\bigr)\,\mathrm d\alpha,
\]
求和即得
\[
\sum_{i=1}^d\mathrm{IG}_i(x)=\int_0^1\frac{\mathrm d}{\mathrm d\alpha}f\bigl(b+\alpha(x-b)\bigr)\,\mathrm d\alpha=f(x)-f(b),
\]
完备性由构造保证，不需要额外公理。它也自然处理了上面那个饱和情形：当前点梯度为零，路径上早先的梯度仍被累积进来。

账单在两处。基线 \(b\) 代表``没有信息时的输出''，可全黑图、全白图、噪声图或训练集均值各是一种诠释，归因图案可以差别很大，而没有哪一个是标准答案。直线路径也不保证落在数据支持集附近，它可能穿过大片没有真实样本的区域，那里的梯度照样计入总账。恢复出来的仍是模型在这条路径上的行为，不因为完备性就升级成因果贡献——改变这个特征、模型输出会怎么变，是另一个需要干预才能回答的问题。

\subsection{解释本身需要被检验}

方法各自内部自洽，不代表结果可用，因而需要一组独立于定义的检验。忠实度问的是解释能不能预测受控扰动下的输出变化：若它说特征 A 最重要，那么单独改动 A 时输出的变化应当与之相符。稳定性问的是输入的微小扰动或随机种子的更换会不会让排名剧烈变化——同一张图连跑三次 Shapley 得到三组不同的前三名，这个解释就不能作为模型的稳定描述。敏感性问的是解释会不会响应：一个永远宣称所有特征等权重要的方法在任何输入上都稳定，也不携带任何信息。适用范围问的是结论针对单个样本、局部邻域还是总体——Shapley 值可以在样本上取平均得到全局重要性，LIME 只在目标点邻域内有效，把一次 LIME 结果当作全局结论是范围误用。

最常用的检验是删除高归因特征、看性能掉多少，这里埋着一个陷阱。删除会制造训练时从未出现过的残缺输入，性能下降里有多少来自特征本身的作用、有多少来自模型在陌生输入上的任意行为，无法从结果里分开。缓解办法是让扰动尽量满足数据生成约束（用条件分布重采样而非置零），并配一份生成机制已知的合成数据作为对照。

\subsection{没有统一正确的解释，只有写清楚的问题}

Shapley 值、LIME 与积分梯度描述的都是给定模型的输入依赖：这个模型在做这次预测时用到了哪些输入信息。它们不回答改变现实世界中的特征会导致什么——那需要干预假设与识别策略，见\hyperref[sec:13-Machine-Learning/12-Causal-Inference/02]{``结构因果模型把干预写成机制替换''}。它们也不回答模型内部是不是在用人类赋予热力图的那个语义——那需要检验内部激活、回路行为，或在表示空间里做反事实干预，而且这类陈述必须通过\hyperref[sec:14-Deep-Learning/10-Interpretability-Mathematics/01]{``神经网络的函数空间视角''}里那道筛子：在参数的等价变换下不变，才谈得上描述模型。

把归因图当作因果证据或机制证据，是范畴错误，哪怕它恰好与直觉吻合——吻合可能说明模型学到了合理的输入依赖，也可能说明任务太简单，很多不同的依赖方式都能给出相似预测。

所以在选方法之前，先把四件事写下来：被解释的是哪个输出量；参照对象是什么；允许对输入做哪类扰动；结论要交给谁去做什么决定。四件事定下来，方法基本也就定了，不同方法之间的差异也从``谁对谁错''变成``问的不是同一件事''。这是本单元反复出现的同一个约束：深度模型把选择表示交给了优化，可验证的保证随之退化为经验规律，剩下能守住的只有一条——把陈述的适用范围写在陈述里。
